\documentclass[11pt]{article}
% Shared declarations; root documents always import a concrete semantic module.
\usepackage{business-admin}
\usepackage{amsmath,amssymb}
\usepackage{longtable,tabularx,array,booktabs}
\usepackage{enumitem}
\usepackage{ragged2e}
\newcolumntype{Y}{>{\RaggedRight\arraybackslash}X}
\newcommand{\SuiteMeta}[2]{\def\DocTitle{#1}\def\DocSubtitle{#2}\def\DocVersion{1.0}\def\DocStatus{Educational reference}\def\DocOwner{Small Business Management Suite}\def\DocSummary{Original teaching and practice resource; no real-business approval implied.}}
\newcommand{\SuiteOpen}{\makedoctitle\vspace{-0.7em}\noindent\rule{\linewidth}{0.45pt}\par\medskip}
\newcommand{\SourceNote}[2]{\par\medskip\noindent{\footnotesize\textbf{Reading anchor:} Longenecker, Petty, Palich and Hoy, \emph{Small Business Management: Launching \& Growing Entrepreneurial Ventures}, 20th ed., Ch.~#1, textbook p.~#2. This document is independently authored commentary and exercises, not a reproduction of the chapter.}}
\newcommand{\WritingLines}[1]{\par\smallskip\noindent\rule{\linewidth}{0.35pt}\par\vspace{0.35cm}\noindent\rule{\linewidth}{0.35pt}\par\vspace{0.35cm}\noindent\rule{\linewidth}{0.35pt}\par\medskip}
\newcommand{\FormField}[1]{\par\medskip\noindent\textbf{#1}\par\WritingLines{3}}
\newcommand{\SBMCheckBox}{\(\square\)}
\newcommand{\AssessmentRow}[2]{#1 & #2 & \(\square\) & \(\square\) & \(\square\) & \(\square\) & \(\square\)\\\hline}
\newcommand{\AssessmentHeader}{\noindent\begin{tabular}{|p{0.22\linewidth}|p{0.49\linewidth}|c|c|c|c|c|}\hline\textbf{Criterion} & \textbf{Evidence to inspect} & \textbf{0} & \textbf{1} & \textbf{2} & \textbf{3} & \textbf{4}\\\hline}
\newcommand{\AssessmentEnd}{\end{tabular}\par}
\setlength{\emergencystretch}{3em}

\SuiteMeta{Case Study: Boutique E-Commerce Launch}{Original fictional cross-chapter simulation}
\begin{document}
\SuiteOpen
\section{Situation}
A direct-to-consumer retailer is considering product expansion while its inventory turns slowly and marketing relies on a single paid channel.
\section{Case evidence (fictional)}
\begin{itemize}[leftmargin=*,itemsep=3pt]
\item A product sells for 64 and costs 26 to acquire and fulfill
\item Per-order promotional acquisition spend averages 18
\item Monthly overhead is 6400
\item The proposed new supplier requires a minimum order of 800 items
\end{itemize}

\section{Student tasks}
\begin{enumerate}[leftmargin=*,itemsep=5pt]
\item Compute contribution after acquisition and a break-even estimate.
\item Model the up-front cash tied in 800 new units.
\item Compare inventory risk with the claimed benefit of broader selection.
\item Design a measured retention test and a supplier due-diligence plan.
\end{enumerate}

\section{Required evidence packet}
\begin{itemize}[leftmargin=*,itemsep=3pt]
\item Decision statement with clearly stated constraints and alternatives.
\item Calculation sheet where numerical claims are made; show units and sensitivity.
\item Risk register with owner, response and residual uncertainty.
\item Recommendation, go/no-go threshold and 30-day follow-up plan.
\end{itemize}

\section{Instructor review notes}
A successful answer does not depend on the same choice as an instructor example. Judge the quality of evidence, arithmetic, critical assumptions, operational feasibility and treatment of adverse effects. Withhold approval when missing financial or compliance information could reverse the recommendation.
\section{Chapter crosswalk}
Review chapters 7, 9, 11, 14, 15, 16, 17, 22. Cases and scenario figures are newly authored; they are not taken from textbook case studies.
\end{document}
